%----------------------------------------------------------------------------------------
% METADATA BLOCK (YAML)  — see inbox/README.md
%----------------------------------------------------------------------------------------
% uid: ""
% school: "lycee"
% level: "2eme"
% section: ["Sciences", "Informatique"]
% chapter: "Notion de Polynômes"
% title: "Produit de quatre nombres consécutifs"
% difficulty: "Moyen"
% tags: ["polynomes", "changement-de-variable", "identites-remarquables", "racine-carree"]
% provenance: "own"
% batch: "2026-11-19-2sc-polynomes-barycentre"
%----------------------------------------------------------------------------------------

\newcommand{\R}{\mathbb{R}}
\newcommand{\vect}[1]{\overrightarrow{#1}}

% ---STATEMENT---
Soit $P(x)=x(x+1)(x+2)(x+3)+1$.
\begin{enumerate}
  \item Développer $x(x+3)$ et $(x+1)(x+2)$.
  \item En posant $t=x^2+3x$, montrer que, pour tout réel $x$ : $P(x)=\left(x^2+3x+1\right)^2$.
  \item En déduire que $\sqrt{97\times98\times99\times100+1}$ est un entier et le calculer.
  \item Résoudre dans $\R$ l'équation $x(x+1)(x+2)(x+3)=24$.
\end{enumerate}

% ---SOLUTION---
\begin{enumerate}[itemsep=14pt]
  \item $x(x+3)=x^2+3x$ \quad et \quad $(x+1)(x+2)=x^2+3x+2$.

  \item On regroupe les facteurs deux par deux (le premier avec le dernier) :
  \[x(x+1)(x+2)(x+3)=\left[x(x+3)\right]\times\left[(x+1)(x+2)\right]=t(t+2)\]
  \[P(x)=t(t+2)+1=t^2+2t+1\]
  \begin{prop}[Identité remarquable]
  \ib{$t^2+2t+1=(t+1)^2$}.
  \end{prop}
  \[P(x)=(t+1)^2=\left(x^2+3x+1\right)^2\]

  \item On prend $x=97$ : $97\times98\times99\times100+1=P(97)=\left(97^2+3\times97+1\right)^2$.

  Comme $97^2+3\times97+1>0$ :
  \[\sqrt{97\times98\times99\times100+1}=97^2+3\times97+1=9\,409+291+1=9\,701\]

  \item $x(x+1)(x+2)(x+3)=24\iff P(x)=25\iff\left(x^2+3x+1\right)^2=25$
  \[\iff x^2+3x+1=5\quad\text{ou}\quad x^2+3x+1=-5\]
  \textbf{Premier cas :} $x^2+3x-4=0$, $\Delta=9+16=25$ :
  \[x=\dfrac{-3-5}{2}=-4\quad\text{ou}\quad x=\dfrac{-3+5}{2}=1\]
  \textbf{Second cas :} $x^2+3x+6=0$, $\Delta=9-24=-15<0$ : pas de solution.
  \[\mathcal S=\{-4\,;1\}\]
  \textbf{Contrôle :} $1\times2\times3\times4=24$ et $(-4)(-3)(-2)(-1)=24$ (vérifié)
\end{enumerate}
\textbf{Idée :} associer les facteurs extrêmes et les facteurs du milieu fait apparaître la même expression $x^2+3x$.

\textbf{Piège :} $A^2=25$ donne \emph{deux} cas, $A=5$ ou $A=-5$.
